\documentclass[letterpaper,twocolumn,10pt]{article}
\usepackage{usenix}
\usepackage{booktabs}

\begin{document}
\date{}
\title{\Large \bf Measuring Semantic Preservation in Passkey Migration}
\author{Anonymous Authors}
\maketitle

\begin{abstract}
An imported passkey can still authenticate while a property needed by an
application has changed. We present PasskeyTransit, a reproducible harness
that measures identity, WebAuthn behavior, extension state, and operational
guarantees separately. A balanced synthetic campaign exercised 6,144
credential--route--repetition attempts; 5,120 reached import and 1,024 were rejected
before a ceremony. Every executed assertion authenticated, yet 512 imported
attempts failed an executed browser behavior oracle for user verification or
largeBlob. These fractions describe the constructed matrix, not product
prevalence. Follow-up controls distinguish a project-authored CXP binding
from selected provider behavior. A Chromium virtual-authenticator transfer
preserved the signing credential but lost its PRF result; four public
assertions from a replay were independently reverified. Fenko authenticated
after CXF import while losing the original HMAC seeds. A Proton software
authenticator preserved both seeds and reproduced the expected PRF with user
verification. We provide claim boundaries and public verifiers. We did not
establish public CXP wire interoperability between providers or a verified
Android transfer.
\end{abstract}

\section{Introduction}
An accepted import and successful login are weak tests of passkey migration.
A relying party can also depend on the credential identifier, WebAuthn
extension output, user-verification behavior, and whether transfer retries
leave duplicate state. The Credential Exchange Format (CXF) represents
credentials~\cite{cxf}; the draft Credential Exchange Protocol (CXP)
addresses protected exchange~\cite{cxp}. WebAuthn authenticates an
assertion~\cite{webauthn}, but a valid signature does not by itself show
that every credential property survived the move.

We model a credential as $C=(I,K,F,S)$: identity, signing material,
observable functions, and operational guarantees. Our falsifiable
proposition is that valid authentication does not imply preservation of
every component. We contribute: (1) a claim matrix and property-specific
oracles; (2) a synthetic campaign with explicit denominators and
non-evaluable properties; (3) executable controls using two project-authored
runtimes, Chromium virtual authenticators, and selected open-source provider
code; and (4) an artifact whose public reports allow independent checks of
retained signatures and PRF digest relationships. This is a measurement
method with selected software cases, not a survey of deployed products.

\section{Background and related work}
CXF can encode FIDO2 HMAC state~\cite{cxf}. The CXP working draft defines a
transport profile~\cite{cxp}; our single-document CXP/HPKE binding is an
experimental project implementation and does not cover its full ZIP/JWE
wire format. HPKE follows RFC~9180~\cite{hpke}. A WebAuthn PRF result is
returned through client extension results, outside the ordinary signed
assertion bytes~\cite{webauthn}. Consequently, signature and PRF oracles
answer separate questions. Ramat et al. compare passkey user experiences
across 111 websites~\cite{ramat}; Jannett et al. measure deployment and
security on passkey-enabled websites~\cite{jannett}. B\"uttner and
Gruschka compare the security and availability of device-bound and synced
credentials~\cite{buttner}. Those units are websites or credential models.
Our unit is an attempted migration of a controlled credential, with
separate post-import oracles for signing, extension behavior, and retry
state. We do not infer website adoption or provider prevalence from it.

\section{Study design}
\subsection{Units, outcomes, and exclusions}
The v1.5 protocol crossed 256 synthetic ES256 credentials (32 in each of
eight strata) with 12 routes and two
repetitions: 6,144 attempts in 96 route-by-stratum groups. The strata
covered basic credentials, PRF with and without UV, largeBlob, credBlob,
a payment-format marker, combined extensions, and a future optional
member. They were constructed controls, not random draws. A route
could reject the document before WebAuthn. For each applicable property we
record PASS, FAIL, NOT\_APPLICABLE, or NOT\_EVALUABLE. The last outcome is
needed where CDP cannot inject PRF or credBlob state.

One oracle checks CXF representation and policy. Each imported attempt
performs two WebAuthn assertions: the first binds a challenge to the attempt,
credential, route, repetition, public key, user handle, RP ID, and origin;
the second binds its challenge to the first and to a canonical digest of
extension observations. A separate verifier checks both ES256 signatures,
the RP hash, user-presence and verification flags, and counter. Extension
oracles examine executed behavior where the
interface permits it; representation loss alone is never relabeled as an
executed PRF failure. A deterministic in-memory simulation tests transfer
atomicity and retry behavior. These oracles are not interchangeable.
Campaign percentages are design-weighted proportions, without population
confidence intervals; paired repetitions are stability checks, not
independent providers. The browser and provider
follow-ups were designed after inspecting earlier results and are
exploratory rather than preregistered confirmations.

\subsection{Threat and implementation boundary}
Two separately written project implementations exchange one CXF document
in Python-to-Node and Node-to-Python directions and reject a modified
ciphertext. This tests that experimental binding. Exporter identity,
authorization, durable provider storage, and the public CXP ZIP/JWE
envelope are outside the tested profile. The transaction model is in
memory and makes no claim about crash durability.

\section{Synthetic campaign results}
Of 6,144 attempts, 1,024 were rejected before ceremony and 5,120 were
imported. Every executed assertion among the imported attempts was
accepted. Table~\ref{tab:c1} separates those accepted assertions from other
oracles. The complete observable PASS yield was 2,304/6,144 (37.50\%).
This describes the balanced test design, not the probability that an
arbitrary user or product will succeed.

\begin{table}[t]
\centering
\caption{C1 synthetic attempts; denominator 6,144.}
\label{tab:c1}
\begin{tabular}{lrr}
\toprule
Classification & Count & Share \\
\midrule
Complete observable PASS & 2,304 & 37.50\% \\
Visible degradation & 960 & 15.62\% \\
Silent degradation & 1,344 & 21.88\% \\
Not evaluable & 512 & 8.33\% \\
Pre-ceremony rejection & 1,024 & 16.67\% \\
\bottomrule
\end{tabular}
\end{table}

Among imported attempts, 512 (10.00\%) combined valid login with failure
of an executed UV or largeBlob browser oracle. PRF or credBlob
representation loss occurred in 1,792 (35.00\%); this category overlaps
the former and does not establish an executed PRF failure in those cells.
The union of non-payment failures was 2,048/5,120 (40.00\%). A
payment-format marker raises the design sensitivity to 2,304/5,120
(45.00\%), but is not a Secure Payment Confirmation ceremony. C2 malformed
document controls comprised 80 cases across ten families and eight strata;
64 were rejected by structural validation, duplicate policy, or strict
preservation. C3 exercised 960 in-memory failure sequences and 1,920 events:
832 rolled back and converged to one copy on retry, while the 128 deliberate
legacy late-failure controls retained a provisional and duplicated on retry.
Neither C2 nor C3 is included in the C1 percentages.

\section{Executable follow-up controls}
\subsection{Experimental exchange and browser PRF loss}
Python and Node recovered identical CXF documents in both directions and
rejected ciphertext alteration. Software PRF evaluators agreed on retained
CXF HMAC seeds and detected their removal. This is a functional software
control, not a browser ceremony.

A Chromium virtual authenticator created a PRF-enabled passkey. Three
source assertions showed a PRF result stable for one salt and different
for a second. The runner exported the signing credential via CDP, carried
the CXF document through the experimental exchange, and injected the
credential into a new virtual authenticator. A fourth assertion
authenticated at the destination, while the PRF result for the original
salt was absent. CDP did not expose a PRF seed. This is a loss control for
that test path, not a finding about a commercial credential provider.

A Windows 2025 GitHub Actions replay recorded six exchange and nine
browser checks as true. A second replay retained four public signatures,
the registration public key, challenges, authenticator and client data,
and SHA-256 digests of observed PRF outputs. A separate verifier
checked four ES256 signatures, RP binding, challenge uniqueness, UV
flags, credential identity, and four PRF digest relationships. The archive
contains no private key, seed, or raw PRF output. Comparing hashes proves
relations among observed outputs, not the HMAC computation from a seed
that CDP did not reveal.

\subsection{Selected provider case: Fenko}
We selected Fenko Vault 0.12.0~\cite{fenko} after inspecting its CXF
path. Two fresh profiles per repetition imported synthetic basic and
PRF-bearing credentials. The credentials signed valid assertions. A
public verifier rechecked 16 signatures and six PRF observation bindings.
The source CXF HMAC seeds were not retained in the observed import/export
path; the browser PRF result differed from an oracle derived from those
source seeds.

A same-provider round trip between two Fenko profiles produced matching
PRF outputs in two selected repetitions. A separate verifier recalculated
24 signatures and 12 PRF bindings. This continuity does not imply that
the original seeds survived: this code path derived a value from the
imported signing key. Functional equality and preservation of the
represented secret are separate claims. Two repetitions are not a
product-level reliability estimate.

\subsection{Positive software control: Proton}
The Proton Pass common 2.0.0 software authenticator~\cite{proton}
provided a selected positive control. In two exploratory repetitions,
the runner imported CXF, exported it, and reimported into a second
software authenticator. Both HMAC seeds survived. With user verification,
the expected PRF was reproduced for two salts before and after reimport.
A modified-seed control changed the result while the credential still
signed. The public verifier rechecked 12 signatures and 12 PRF digest
comparisons, including four modified-seed controls. The isolated report
cannot reconstruct the HMAC from a withheld seed; that requires repeating
the pinned software run. This route used pinned WASM, not a
browser passkey provider or the installed Proton Pass application; no
no-UV ceremony was run.

\section{Discussion and limitations}
The browser loss control and Fenko case show different ways a signing
credential can outlive its original PRF state. The Proton software route
shows that preserving CXF seeds and reproducing the expected PRF is
possible in one selected implementation. These cases are not a
representative sample of provider behavior.

The public CXP wire format was not exchanged between two independent
providers. A human-reported Android Bitwarden-to-Google transfer lacks an
instrumented destination assertion and transfer receipt and is excluded
from measured results. CDP cannot inject original PRF state; positive
browser-to-browser restoration of CXF HMAC seeds is NOT\_EVALUABLE. The
Proton module result cannot fill that gap. The transaction simulation
cannot support claims about storage durability, and the synthetic campaign
cannot support product prevalence or superiority claims.

\section{Conclusion}
Passkey migration needs separate oracles for identity, signing,
extension behavior, and transaction semantics. PasskeyTransit records
when an interface makes a property unobservable. Its designed campaign
and selected software controls demonstrate that valid authentication can
coexist with a changed or missing property. The public replay lets readers
recompute browser-control signatures without private credential material.
Provider interoperability and verified mobile transfer require separate
experiments.

\appendix
\section{Open Science}
The review package contains source, locked dependencies, protocols,
synthetic reports, two GitHub Actions ZIP archives for the browser
control, public signature transcripts, and verification scripts.
\texttt{python tools/verify\_v08\_draft.py} checks pinned hashes and the
phase-14, Fenko, and Proton public reports. A clean extraction with newly
installed locked dependencies passed the same checks on 20 October 2026
on the same team's Windows host. This was an environment-isolation check,
not independent replication.
Provider vaults, private keys, original PRF seeds, and raw PRF outputs
are excluded. Repeating provider paths requires pinned external revisions.
An anonymous artifact URL without tracking must be provided before
submission; the current public repository is identifying.

\section{Ethical considerations}
All generated credentials and RP accounts in the reported software
experiments were synthetic. The browser control used virtual
authenticators and a localhost RP. No real user passkey or private vault
is distributed. The excluded Android observation is not presented as
independently verified evidence. We report provider behavior with route
and revision scope rather than a general vendor security defect.

\begin{thebibliography}{16}
\bibitem{cxf} FIDO Alliance. Credential Exchange Format v1.0 Proposed Standard Errata, 2026. \url{https://fidoalliance.org/specs/cx/cxf-v1.0-ps-errata-20260309.html}.
\bibitem{cxp} FIDO Alliance. Credential Exchange Protocol v1.0 Working Draft, 3 October 2024. \url{https://fidoalliance.org/specs/cx/cxp-v1.0-wd-20241003.html}.
\bibitem{hpke} R. Barnes et al. Hybrid Public Key Encryption. RFC 9180, 2022. \url{https://www.rfc-editor.org/rfc/rfc9180}.
\bibitem{webauthn} W3C. Web Authentication Level 3, Recommendation, 25 August 2026. \url{https://www.w3.org/TR/2026/REC-webauthn-3-20260825/}.
\bibitem{ramat} B. Ramat, D. Kartchner, M. Clark, and K. Seamons. Passkeys in the Wild: A Systematic Study of FIDO2 User Experience Consistency Across Websites. SOUPS, pp. 21--39, 2026. \url{https://www.usenix.org/conference/soups2026/presentation/ramat}.
\bibitem{jannett} L. Jannett, A. Mayer, M. Westers, V. Mladenov, C. Mainka, and J. Schwenk. The State of Passkeys: Studying the Adoption and Security of Passkeys on the Web. USENIX Security, pp. 5039--5058, 2026. \url{https://www.usenix.org/conference/usenixsecurity26/presentation/jannett}.
\bibitem{buttner} A. B{\"u}ttner and N. Gruschka. Device-Bound vs. Synced Credentials: A Comparative Evaluation of Passkey Authentication. ICISSP, 2025. \url{https://arxiv.org/abs/2501.07380}.
\bibitem{fenko} FenkoHQ. Passkey Vault 0.12.0, pinned commit \texttt{2d873fbe25257250de6e322e9d785b97b17a544f}. \url{https://github.com/FenkoHQ/passkey-vault}.
\bibitem{proton} Proton. Proton Pass common 2.0.0, pinned commit \texttt{69a0927f0fe7071243d6c38e404459a69210fdd4}. \url{https://github.com/protonpass/proton-pass-common}.
\end{thebibliography}
\end{document}
